\section{Introduction}\label{sec:introduction}

Grid operators required robust day-ahead insights into where load will increase and how strongly the system is bound. In developing economies, accelerated demand growth commonly corresponds with supply limits and load shedding, so planners cannot depend on demand forecasting alone; they also need a understandable signal of operational stress. Bangladesh’s national power system reports daily data for nine administrative divisions, with evening-peak demand emerging together with reserve, limitation, and shedding fields in the same records \cite{Hossain2025BangladeshSmartGrid}. Planning is therefore fundamentally spatial and system-wide at once.

In practice, however, load forecasting and stress monitoring are still widely developed as individual tools. Traditional and combined models \cite{Breiman2001RandomForest,Chen2016XGBoost,Musbah2025LoadForecasting} predict demand without modelling connections between divisions. Shallow recurrent or graph-convolution networks \cite{Hochreiter1997LSTM,Zhao2020TGCN} may underperform to merge a full weekly lookback with explanatory cross-node framework. Graph methods established for residential or microgrid settings often presume geographical neighbors \cite{Kipf2017GCN}, which may not align how regional loads in effectively move together on a countrywide footprint. Operational stress- influenced by limitation pressure, reserve shortfall, and shedding- is rarely forecast collaboratively with regional demand; most shedding work \cite{Paphitis2025SpatialLoadShedding} targets control instead of next-day stress indices.  \cite{Baur2024XAIReview,Lundberg2017SHAP}.

We treat day-ahead forecasting as a supervised multi-task problem \cite{Caruana1997Multitask} over nine divisions. Task~1 predicts next-day regional evening-peak demand in MW. Task~2 predicts a graph-level Operational Stress Index (OSI) in $[0,1]$ that summarises shedding intensity, reserve margin, and limitation pressure. Demand and OSI are related in operations but not redundant: demand describes expected load by region, while OSI captures constraint pressure that demand alone does not fully explain. The question we ask is whether one spatio-temporal model can learn both outputs from historical Bangladesh grid records under strict chronological evaluation.

Our approach is the Correlation-Aware Multi-Task Forecasting Framework, implemented as a Parallel-Fusion Spatio-Temporal Graph Transformer (PF-STGT). The model takes a seven-day lookback of division-level and national features together with a correlation graph estimated from training data, combines graph-aware spatial attention with temporal transformers \cite{Vaswani2017Attention}, and decodes regional demand and bounded OSI through separate heads. Training uses a repaired multi-task loss so neither task is drowned out during optimisation. We evaluate against classical, recurrent, and graph-convolution baselines, run configuration ablations, apply Wilcoxon testing with Bonferroni correction, and report post-hoc attributions for the selected model.

This paper makes five contributions: (i)~a correlation-graph PF-STGT for joint day-ahead demand and OSI forecasting on Bangladesh national data; (ii)~evidence that correlation-only adjacency outperforms geographical and hybrid priors on this dataset; (iii)~a practical multi-task loss calibration that stabilises joint training; (iv)~a full benchmark and ablation programme with inferential testing; and (v)~integrated explainability with multi-method attribution for both tasks.

Section~\ref{sec:methodology} describes the method; Sections~\ref{sec:experimental_setup}--\ref{sec:results} present experiments and results; Sections~\ref{sec:discussion}--\ref{sec:conclusion} interpret the findings and conclude.
